\documentclass[11pt]{amsart}

\usepackage[T1]{fontenc}
\usepackage{lmodern}
\usepackage{microtype}
\usepackage{mathtools,amssymb,amsthm}
\usepackage[hidelinks]{hyperref}

\hypersetup{
  pdftitle={Short-Cycle Counterexamples to Conjecture 12 of Berman--Gevay--Pisanski}
}

\newtheorem{theorem}{Theorem}
\newtheorem{corollary}[theorem]{Corollary}
\theoremstyle{remark}
\newtheorem{remark}[theorem]{Remark}

\newcommand{\Z}{\mathbb Z}
\newcommand{\GammaOne}{\Gamma_1(m;a,b)}
\newcommand{\GammaTwo}{\Gamma_2(m;a,b)}

\title[Short-cycle counterexamples]
{Short-Cycle Counterexamples to Conjecture~12 of
Berman--G\'evay--Pisanski}

\date{}

\subjclass[2020]{Primary 05B30; Secondary 51E30, 05C25, 52C30}
\keywords{polycyclic configuration, voltage graph, Levi graph, girth}

\begin{document}

\begin{abstract}
Berman, G\'evay, and Pisanski proposed two parameterized cyclic-voltage
families for geometric $(7m)_4$ configurations, while noting that further
parameter restrictions might be necessary. We determine exactly when the
corresponding derived multigraphs avoid cycles of length two and four. This
refutes the sufficiency of the parameter ranges displayed in their
Conjecture~12: the
derived graph of $F_2(m;2,1)$ has a $4$-cycle for every $m\ge5$, and that of
$F_1(m;m/2,b)$ has a $4$-cycle for every even $m\ge4$ and every
$1\le b\le m/2$. Hence these assignments do not yield even combinatorial
configurations. We do not address geometric realizability of the surviving
lifts.
\end{abstract}

\maketitle

\section{The voltage families}

For $m>3$, Berman, G\'evay, and Pisanski introduced a generic reduced Levi
graph $H=\operatorname{RLG}(B)$ and, in Conjecture~12 of~\cite{BGP}
(Conjecture~7.1 of the preprint version listed there, whose statement is
identical), proposed two families of cyclic voltage assignments over $\Z_m$,
in the parameter order
\[
(a,c,d,e,f,g,q,a',c',d',e',f',g',q',t):
\]
\begin{align*}
F_1(m;a,b)
 &= (a,-a,a,b,b,-a,-2a,
     a,-a,a,b,b,-a,-2a,0),\\
F_2(m;a,b)
 &= (a,b,b,b,b,b,b,
     a,b,b,b,b,b,b,0).
\end{align*}
The displayed range for $F_1$ is
\[
a\ge b,\qquad 1\le a,b\le m/2,
\]
and that for $F_2$ is
\[
a\ge b,\qquad a\ne m/2.
\]
The conjecture is stated with the proviso that ``there may be other
constraints on $a$ and $b$ that have not yet been identified''~\cite{BGP}.
Our purpose is to determine exactly the constraints forced by cycles of
length at most four.

The line-orbit vertices of $H$ are
\[
M,R,P,B,C,Y,G,
\]
and the point-orbit vertices are
\[
r,m,p,y,g,c,b.
\]
This notation is inherited from~\cite{BGP} and overloads four letters: $b$,
$c$ and $g$ name both point orbits and voltages, and $m$ is both the modulus
and a point orbit. Subscripted lower-case letters always denote lifted
points, so $m_1$ in the proof of Corollary~\ref{cor:counterexamples} is a
point of the derived graph and not the modulus.

Using the standard cyclic-voltage formalism for polycyclic
configurations \cite{BP}, and following~\cite[Sec.~2]{BGP}, an edge of
voltage $\alpha\in\Z_m$ means that the lifted line $L_i$ is incident with
the lifted point $v_{i+\alpha}$. Transcribing the generic assignment drawn in
Figure~4(b) of~\cite{BGP} into incidence form gives
\begin{equation}\label{eq:incidences}
\begin{aligned}
M&:\ r(0),r(a),y(c),p(q),
& R&:\ m(0),m(a'),g(d),p(0),\\
P&:\ r(0),y(0),g(0),m(q'),
& B&:\ y(e'),g(0),c(g),c(t),\\
C&:\ y(0),g(f'),b(0),b(g'),
& Y&:\ m(c'),b(e),c(0),p(0),\\
G&:\ r(d'),b(0),c(f),p(0).
\end{aligned}
\end{equation}

\medskip
\noindent\textbf{Provenance of \eqref{eq:incidences}.}
Figure~4(b) of~\cite{BGP} is a drawing, so \eqref{eq:incidences} is a
transcription of it rather than a quotation, and everything below is a
statement about the graph that \eqref{eq:incidences} specifies. We therefore
print the transcription in full and record the checks that pin it down. Of
the two checks that follow, the structural one is internal to
\eqref{eq:incidences} and can be carried out from the displayed data alone;
the numerical one is the single point at which this note appeals to data
outside it, namely the incidence table of~\cite[Table~1]{BGP}, which we do
not reproduce here.

Structurally, \eqref{eq:incidences} lists $28$ edges on the $14$ vertices of
$H$, four at each vertex; the $13$ edges of voltage $0$ form a spanning tree
of $H$, and the $15$ named voltages occupy the $15$ remaining edges, one
each. This is the normalization realized by the fifteen edge labels
$a,c,d,e,f,g,q,a',c',d',e',f',g',q',t$ carried by the figure.

Numerically, substituting $F_1(3;1,1)$, that is
\begin{multline*}
(a,c,d,e,f,g,q,a',c',d',e',f',g',q',t)\\
=(1,2,1,1,1,2,1,1,2,1,1,1,2,1,0)\quad\text{in }\Z_3,
\end{multline*}
into \eqref{eq:incidences} reproduces all $28$ incidences of the incidence
table of $B(21_4)$ in~\cite[Table~1]{BGP} after the relabelling
$i\mapsto-i$ of $\Z_3$, and \cite{BGP} record that $F_1(3;1,1)$ is isomorphic
to $B(21_4)$. The relabelling is harmless, because $i\mapsto-i$
is an automorphism of the derived graph and so preserves girth. In particular
the four parallel pairs of
\eqref{eq:incidences}, namely $Mr=\{0,a\}$, $Rm=\{0,a'\}$, $Bc=\{g,t\}$ and
$Cb=\{0,g'\}$, then carry the voltage differences $1,1,2,2$, matching the
four points of that table which lie on two lines of a single orbit: $r_0$ on
$M_0,M_1$; $m_0$ on $R_0,R_1$; $c_0$ on $B_0,B_2$; $b_0$ on $C_0,C_2$.

At $m=3$ with $a=b=1$ several of the fifteen voltages take equal values, so
this comparison alone cannot detect a permutation of labels within a block of
equal values. The degeneracy is genuine: at $F_1(3;1,1)$ the generic values
$a$, $b$ and $-2a$ all equal $1$ in $\Z_3$, so the comparison of parallel
pairs shows only that the label carried by $Mr$ has generic value $a$, $b$ or
$-2a$, and not that it is $a$ itself. We therefore do not claim that
\eqref{eq:incidences} is the only reading of Figure~4(b) compatible with these
checks, and we claim nothing beyond it: Theorem~\ref{thm:classification} and
Corollary~\ref{cor:counterexamples} are assertions about the graph that
\eqref{eq:incidences} specifies. Consequently a reader without
\cite[Table~1]{BGP} at hand can still check everything in the sections that
follow, taking \eqref{eq:incidences} as the definition of $H$ and of its
voltages; what rests on the comparison with that table is only the
identification of this graph with the one drawn in Figure~4(b)
of~\cite{BGP}, and hence the bearing of our results on Conjecture~12 as its
authors intended it. The four parallel pairs and the ten
quadrilaterals used in the proof below are features of $H$ itself, so a reader
who reads the figure differently need only substitute that reading into
\eqref{eq:incidences} and repeat the proof with the labels permuted
accordingly.

\medskip
Let $\Gamma_i(m;a,b)$ denote the derived multigraph associated with
$F_i(m;a,b)$. A Levi graph of a combinatorial point-line configuration is
simple and has no $4$-cycle; thus the relevant derived graph must have
girth at least six, where a parallel pair is regarded as a cycle of length
two.

\section{Short-cycle classification}

\begin{theorem}\label{thm:classification}
Let $m\ge2$ and let $a,b\in\Z_m$.
\begin{enumerate}
\item
The graph $\GammaOne$ has girth at least six if and only if
\[
a\ne0,\qquad b\ne0,\qquad 2a\ne0,\qquad 2b\ne0
\quad\text{in }\Z_m.
\]
\item
The graph $\GammaTwo$ has girth at least six if and only if
\[
a,b,2a,2b,a-b,a-2b\ne0
\quad\text{in }\Z_m.
\]
\end{enumerate}
\end{theorem}

\begin{proof}
A cycle in a derived voltage graph projects to a cyclically reduced closed
walk of the same length in the base graph, and the net voltage of that walk
is zero. Conversely, a cyclically reduced zero-voltage closed walk lifts to
a nonbacktracking closed walk and therefore contains a cycle of no greater
length. Since the derived graphs are bipartite, it is enough to enumerate
the projected closed walks of lengths two and four.

The four parallel pairs in $H$ are $Mr$, $Rm$, $Bc$, and $Cb$. Their
voltage differences, up to sign, are
\[
a,\qquad a',\qquad g-t,\qquad g',
\]
and traversing a parallel pair twice gives
\[
2a,\qquad 2a',\qquad 2(g-t),\qquad 2g'.
\]
The four parallel pairs are pairwise vertex-disjoint, so no vertex of $H$
lies in two of them. A cyclically reduced closed walk of length four visits
four, three, or two distinct vertices. On three vertices it has the form
$L-v_1-L-v_2-L$ or $L_1-v-L_2-v-L_1$; the first requires the two pairs
$Lv_1$ and $Lv_2$ to be parallel and the second requires $vL_1$ and $vL_2$ to
be parallel, so each would place the repeated vertex, $L$ respectively $v$,
in two distinct parallel pairs, and both are excluded. On two vertices it
traverses a single parallel pair twice, which is the case already treated.
Hence every remaining cyclically reduced closed walk of length four visits
four distinct base vertices and traces one of the ten quadrilaterals below.
Directly
from~\eqref{eq:incidences}, their net voltages are as follows, up to sign:
\[
\begin{array}{c|c@{\qquad}c|c}
\text{walk} & \text{net voltage} & \text{walk} & \text{net voltage}\\ \hline
MrPyM & c,\ a-c
  & RmYpR & c',\ a'-c'\\
GbYpG & e
  & GcYpG & f\\
GbYcG & e+f
  & ByPgB & e'\\
CyPgC & f'
  & ByCgB & e'+f'\\
GrMpG & d'+q,\ a-d'-q
  & PgRmP & d+q',\ a'-d-q'
\end{array}
\]
These four parallel pairs and ten quadrilaterals exhaust the cyclically
reduced closed walks of lengths two and four in $H$.

Substitution of $F_1$ leaves, up to sign, precisely
\[
a,\qquad b,\qquad 2a,\qquad 2b.
\]
Substitution of $F_2$ leaves, up to sign, precisely
\[
a,\qquad b,\qquad 2a,\qquad 2b,\qquad a-b,\qquad a-2b.
\]
The two assertions follow.
\end{proof}

\begin{corollary}\label{cor:counterexamples}
The displayed parameter ranges in Conjecture~12 of~\cite{BGP} are not
sufficient.
\begin{enumerate}
\item
For every $m\ge5$, the assignment $F_2(m;2,1)$ satisfies the displayed
conditions, and even
\[
1\le1<2<m/2,
\]
but its derived graph has a $4$-cycle. In particular,
\[
F_2(5;2,1)
=(2,1,1,1,1,1,1,2,1,1,1,1,1,1,0)
\]
does not yield a $(35_4)$ configuration, geometric or combinatorial.
\item
For every even $m\ge4$ and every $1\le b\le m/2$, the assignment
$F_1(m;m/2,b)$ satisfies the displayed conditions but its derived graph
has a $4$-cycle.
\end{enumerate}
\end{corollary}

\begin{proof}
For $F_2(m;2,1)$, equation~\eqref{eq:incidences} gives
\[
P_0\sim g_0,m_1,
\qquad
R_{-1}\sim g_0,m_1.
\]
The vertices $P_0,R_{-1},g_0,m_1$ are distinct, so
\[
P_0-g_0-R_{-1}-m_1-P_0
\]
is a $4$-cycle. Equivalently, the obstruction $a-2b$ in
Theorem~\ref{thm:classification} vanishes. For $m\ge5$, the choice
$(a,b)=(2,1)$ satisfies the displayed conditions of Conjecture~12 and also
the stated strict half-range.

Now let $m$ be even and set $a=m/2$. The two lifted lines $M_0$ and $M_a$
both contain the two distinct points $r_0$ and $r_a$, because $2a=0$ in
$\Z_m$. Hence
\[
M_0-r_0-M_a-r_a-M_0
\]
is a $4$-cycle. A point-line configuration cannot have two distinct lines
through two distinct common points, so none of these assignments yields the
asserted configuration.
\end{proof}

\begin{remark}
Conjecture~12 explicitly allows that further restrictions may be necessary.
Accordingly, the result above corrects the displayed sufficiency ranges; it
does not settle geometric realizability for the remaining parameter values.

Since $\Gamma_i(m;a,b)$ is bipartite and $4$-regular with parts of size $7m$,
girth at least six is not merely necessary but equivalent to the lift being a
combinatorial $(7m)_4$ configuration. Theorem~\ref{thm:classification}
therefore determines exactly which of these assignments yield combinatorial
configurations, and geometric realizability is the only question it leaves
open.

Within the printed range for $F_1$, Theorem~\ref{thm:classification} says
that the short-cycle condition is exactly $a<m/2$. The printed range for
$F_2$ places no upper bound on $a$ and does not separate $a$ from $b$, so
besides $a=2b$ it also admits $a=b$, for instance $F_2(m;1,1)$ for every
$m>3$, and $b=m/2$ with
$a>m/2$, for instance $F_2(6;4,3)$; these make $a-b$ and $2b$ vanish
respectively. The $F_2$ range is thus insufficient for every $m>3$, not only
for $m\ge5$. Within the stronger range $1\le b<a<m/2$ it is exactly
$a\ne2b$.

None of this contradicts an assignment displayed in~\cite{BGP}. The
illustrated examples $F_1(4;1,1)$, $F_1(5;2,2)$, $F_1(6;2,1)$ and
$F_2(4;3,1)$, $F_2(5;3,1)$, $F_2(6;4,1)$ of~\cite[Figs.~11 and~12]{BGP} all
satisfy the criteria of Theorem~\ref{thm:classification}: in each of the three
$F_1$ examples none of $a$, $b$, $2a$, $2b$ vanishes, and in each of the three
$F_2$ examples $a-b$ and $a-2b$ do not vanish either. The two lists of
criteria differ, and the $F_1$ example $F_1(6;2,1)$ does have $a=2b$; by
Theorem~\ref{thm:classification} that is an obstruction for $F_2$ and not for
$F_1$. The counterexamples above occupy parameter values that
lie inside the printed ranges but were not illustrated. That $F_2(4;3,1)$ has
$a>m/2$ also confirms that the bound $1\le a,b\le m/2$ is imposed on $F_1$
alone. We are not aware of prior work identifying the missing constraints.
\end{remark}

\begin{thebibliography}{9}

\bibitem{BGP}
L.~W. Berman, G. G\'evay, and T. Pisanski,
\emph{On a new $(21_4)$ polycyclic configuration},
Electron. J. Combin. \textbf{31} (2024), no.~4,
Paper No.~P4.54, 26~pp.,
\href{https://doi.org/10.37236/12405}{doi:10.37236/12405};
preprint version \href{https://arxiv.org/abs/2309.12992}{arXiv:2309.12992},
in which the statement cited here as Conjecture~12 is numbered
Conjecture~7.1.

\bibitem{BP}
M. Boben and T. Pisanski,
\emph{Polycyclic configurations},
European J. Combin. \textbf{24} (2003), no.~4, 431--457,
\href{https://doi.org/10.1016/S0195-6698(03)00031-3}
{doi:10.1016/S0195-6698(03)00031-3}.

\end{thebibliography}

\end{document}